\documentclass[11pt]{article}
\usepackage{fontspec}
\setmainfont{Times New Roman}
\setsansfont{Arial}
\setmonofont{Consolas}
\usepackage{amsmath,amssymb}
\usepackage{geometry}
\usepackage{microtype}
\geometry{a4paper,margin=3cm}
\setlength{\parindent}{1.25em}
\setlength{\parskip}{0pt}
\emergencystretch=30em
\allowdisplaybreaks
\raggedbottom
\providecommand{\Rfield}{\mathfrak R}

\begin{document}

% Унит: Папер 11. Соурце сегментс П11-С0001--П11-С0027.

\begin{center}
{\Large\bfseries 11. Уравненьа с заданој групој}\par
\vspace{1.5em}
\mbox{Math. Ann. 78 (1918), s. 221--229}
\end{center}

\vspace{4em}

Проблем конструкције уравнень с заданој групој можно нападати в двух смерјах, кторе можно кратко назвати «ирационалным», то јест характеризујучим корени, и «рационалным», то јест характеризујучим коефициенты.

К «ирационалному» смерју, кторы работаје функционално-теоретично и аритметично, принадлежи теорем Кронекера, же все Абелове польа в области рационалных чисел сут циклотомичне польа, и одповедајучи теоремы за релативно-Абелове польа в односу к квадратным числовым польам. Ты теоремы дају, с једнеј страны, реализујемост всих Абеловых груп в односу к тым специалным областјам рационалности; с другој страны, они дају целост уравнень и заједно глубши поглед в структуру числовых поль, определьеных теми уравненьами. Всакы поједины теорем једнак јест ограничены на специалну област рационалности; и всака нова област рационалности потребује совсем ново обработанье.

«Рационалны», алгебраичны смер опираје се на Хилбертов теорем о неразложимости, по кторому в представјеньу коефициентов можно ограничити се на параметрно представјенье. Сем принадлежи, напр., обстојанье либовольно многих уравнень с алтернативној групој в односу к либовольно заданој области рационалности. Туто природно недостаје горе означены поглед в структуру поль, дефинираных уравненьами; але предност тој «рационалној» смери лежи в том, же реализујемост груп доказује се за {\itshape все области рационалности једночасно}; однако до тепер познане параметрне представјеньа в обче не дају целост уравнень.\footnote{За разрешиме уравненьа параметрно представјенье коефициентов можно заменити представјеньем коренов. Недавно Ф. Мертенс дал за неке групы 8-го степена најобчејше представјеньа коренов: \emph{Gleichungen 8ten Grades mit Quaternionengruppen}, \emph{Sitzb. d. Ak. d. Wiss., Wien, Abt. IIa, Bd. 125, S. 735}. Како из тој ноты беру, Г. Бухт уже ранеје дал најобчејше изразы коренов за всегда конструователне нормалне формы уравнень 3-го и 4-го степена и уравнень 8-го степена с поменьеными групами: \emph{Über einige algebraische Körper achten Grades}, \emph{Arkiv for Math., Astron. och Physik, Bd. 6, Nr. 30}. [Додаток при коректури.]} 

То, чо следује, јест принос к тој «рационалној» смери: дају се {\itshape достаточне условја} за {\itshape обстојанье параметрных представјень, кторе при либовольној области рационалности дају целост уравнень с заданој групој}. Ты условја стојат в том, же инвариантно полье, принадлежече групе, -- Лагранжев родовы домен, -- можно рационално представити чрез {\itshape независне} функције области, то јест оно имаје «минималну базу»; или, другыми словами, јест изоморфно с польем всих рационалных функциј од $n$ неизвестных. Како буде показано јешче в § 2, то јест, напр., исполньено за все групы, кторе наступајут при уравненьах 3-го и 4-го степена. Стварно поставјенье и точнејша дискусија тых уравнень 3-го и 4-го степена находи се в дисертацији Ф. Сеиделманна,\footnote{\emph{Die Gesamtheit der kubischen und biquadratischen Gleichungen mit Affekt bei beliebigem Rationalitätsbereich}. Ерланген 1916.} из кој безпосредно следује извод тој објавы.

Питанье о минималној бази Лагранжевых родовых доменов было -- како уже поменьено в работе \emph{Körper und Systeme rationaler Funktionen} (\emph{Math. Ann. 76}) -- поставльено мне од Е. Фишера независно од описаној связи и дало побудку к настојачим изследованьем.\footnote{Срав. чест 2 предварительној објавы о \emph{Rationale Funktionenkörper}, \emph{Jahresb. d. d. Math. Vereinig., Bd. 22, 1913}.}

\begin{center}
§ 1.\\[0.75\baselineskip]
\textbf{Достаточне условја за параметрно представјенье.}
\end{center}

\noindent\textbf{1.}\quad Можност свести конструкцију уравнень с заданој групој на параметрно представјенье следује из Хилбертового теорема о неразложимости, кторы говори следече. Нехај уравненье
\begin{equation}
 f(x)=x^n+F_1(\lambda_1\cdots\lambda_r)x^{n-1}+\cdots+F_n(\lambda_1\cdots\lambda_r)=0
\tag{1}
\end{equation}
-- где $F_i(\lambda_1\cdots\lambda_r)$ значет рационалне функције независных параметров $\lambda_1\cdots\lambda_r$, с коефициентами из числового польа\footnote{Вместо числового польа могла бы наступити такоже обча област рационалности.} $\Omega$ -- имаје, в односу к польу $\Omega(\lambda_1\cdots\lambda_r)$ всих рационалных функциј од $\lambda_1\cdots\lambda_r$ с коефициентами из $\Omega$, групу $\Gamma$. Тогда параметры $\lambda$ можно заменити либовольно многими рационалными числами тако, же взникнуше числово уравненье
\begin{equation}
 g(x)=x^n+a_1x^{n-1}+\cdots+a_n=0
\tag{2}
\end{equation}
в односу к $\Omega$ имаје управо групу $\Gamma$. Такоје $g(x)$ будемо точнејше означати чрез $g_\Gamma$.

Тепер пытаје се, чи можно {\itshape дати такоје параметрно представјенье} (1), {\itshape же взникнуше числово уравненье} (2) {\itshape пробегаје целост всих $g_\Gamma$, когда $\lambda_1\cdots\lambda_r$ пробегајут все величины из $\Omega$}, -- с изкльученьем тых величин из $\Omega$, кторе вызывајут редукцију групы.\footnote{Сем треба причислити такоже појавјенье двојнего кореньа уравненьа $g(x)=0$.} Другыми словами: нелинеарны систем уравнень
\begin{equation}
 F_1(\lambda_1\cdots\lambda_r)=a_1;\quad \cdots;\quad F_n(\lambda_1\cdots\lambda_r)=a_n
\tag{3}
\end{equation}
треба имети решенье за всакы систем вредностиј $a_1\cdots a_n$, кторы одповедаје некому $g_\Gamma$, и то в величинах $\lambda_1\cdots\lambda_r$ из $\Omega$. Понеже систем (3) јест нелинеарны, треба од почетка положити ограниченье в пожаданьу, же целост $g_\Gamma$ буде достана чрез параметрно представјенье. Именне, при такых нелинеарных системах могут всегда наступити сингуларне системы вредностиј $a_1\cdots a_n$, кторе удовлетворјајут алгебраичној релацији $H(a_1\cdots a_n)=0$ и за кторе {\itshape не} обстојит решенье;\footnote{То скоро буде обче проведено в работе о «алтернативе при нелинеарных системах уравнень».} тогда параметрно представјенье (1) не успеваје и потребно јест дополнително представјенье. Але в настојачем случају ты сингуларне вредности $a_1\cdots a_n$, како буде показано, можно просто характеризовати.

\noindent\textbf{2.}\quad Да бы добити такоје параметрно представјенье, кладемо силнејше пожаданье, же $\lambda_i$ сут {\itshape природне ирационалности принадлежече групе, идентично в кореньах смотреных како неизвестне}. Под тым разумејемо следече. Ако в (1) заменимо параметры $\lambda_1\cdots\lambda_r$ пригодно избраными {\itshape рационалными} функцијами $\varphi_1(x)\cdots\varphi_r(x)$ неизвестных $x_1\cdots x_n$ -- с коефициентами из $\Omega$ --, тогда (1) имаје прејти в
\begin{equation}
 h(x)=x^n-\sigma_1(x)x^{n-1}+\sigma_2(x)x^{n-2}\cdots\pm\sigma_n(x),
\tag{4}
\end{equation}
-- где $\sigma_1(x)\cdots\sigma_n(x)$ значет елементарне симетричне функције од $x_1\cdots x_n$ --; и $h(x)$ в односу к области рационалности $\Omega(\varphi_1(x)\cdots\varphi_r(x))$, котора при том настаје из $\Omega(\lambda_1\cdots\lambda_r)$, имаје имети групу $\Gamma$. Заради независности параметров функције $\varphi_1(x)\cdots\varphi_r(x)$ такоже имајут быти алгебраично независне.

За прецизованье тој пожаданьа треба објаснити связ $\Omega(\varphi_1(x)\cdots\varphi_r(x))$ с групој $\Gamma$. За то најпрво определимо најобчејшу област рационалности $K$, в односу к кој $h(x)$ ставаје $h_\Gamma$. Нехај $\mathfrak L_\Gamma$ означаје полье всих рационалных функциј с рационално-числовыми коефициентами од $x_1\cdots x_n$, кторе допушчају $\Gamma$ -- такоже называно Лагранжев родовы домен принадлежечи $\Gamma$, или инвариантно полье $\Gamma$; и нехај $\Omega(\mathfrak L_\Gamma)$ означаје одповедајучи полье, настајуче чрез адјункцију $\mathfrak L_\Gamma$ к $\Omega$. Теды $\mathfrak L_\Gamma$ јест подполье польа $\mathfrak R(x_1\cdots x_n)$ всих рационалных функциј с рационално-числовыми коефициентами од $x_1\cdots x_n$. По дефиницији групы најобчејша област $K$ дана јест чрез {\itshape всако полье, кторо садржаје $\mathfrak L_\Gamma$ како подполье и којего пререз с $\mathfrak R(x_1\cdots x_n)$ јест точно $\mathfrak L_\Gamma$}. Одповедајуче, за всаку област $K$, ктора садржаје $\Omega$, пререз с $\Omega(x_1\cdots x_n)$ јест дан чрез $\Omega(\mathfrak L_\Gamma)$. Значи особливо, понеже по нашему пожаданьу $\Omega(\varphi_1(x)\cdots\varphi_r(x))$ имаје быти област $K$, пререз $\Omega(\varphi_1(x)\cdots\varphi_r(x))$ с $\Omega(x_1\cdots x_n)$ дан јест чрез $\Omega(\mathfrak L_\Gamma)$; а понеже с другој страны $\Omega(\varphi_1(x)\cdots\varphi_r(x))$ јест подполье $\Omega(x_1\cdots x_n)$, оно ставаје точно равно $\Omega(\mathfrak L_\Gamma)$. Пожаданье значи теды, же {\itshape Лагранжев родовы домен $\Omega(\mathfrak L_\Gamma)$ имаје настати чрез адјункцију алгебраично независных функциј $\varphi_1(x)\cdots\varphi_r(x)$ к $\Omega$}; при том мора быти $r=n$, понеже $\mathfrak L_\Gamma$ садржаје $n$ алгебраично независных елементарных функциј, а с другој страны не може быти веч неж $n$ алгебраично независных функциј. Функције $\varphi_1(x)\cdots\varphi_n(x)$ будут, како јешче будемо говорити, творити {\itshape минималну базу} $\Omega(\mathfrak L_\Gamma)$; или такоже $\mathfrak L_\Gamma$ имаје минималну базу с коефициентами из $\Omega$.

\noindent\textbf{3.}\quad Тепер суштно јест, же цело то разсмотренье можно обратити; и тако оно стварно води к достаточным условјам\footnote{Пожаданье в ч. 2 было разширјено в односу к ч. 1, тако же условја не морајут быти необходне.} за искано параметрно представјенье. Предполагајемо теды, же $\varphi_1(x)\cdots\varphi_n(x)$ јест минимална база $\Omega(\mathfrak L_\Gamma)$; и нехај $\Omega$ буде најманьше полье, садржече коефициенты $\varphi_1(x)\cdots\varphi_n(x)$. Тогда $\Omega$ необходно јест алгебраично числово полье, и може особливо быти полье $\mathfrak R$ всих рационалных чисел. Понеже $\sigma_1(x),\ldots,\sigma_n(x)$ принадлежет к $\Omega(\mathfrak L_\Gamma)$, обстојат идентичности в $x_1\cdots x_n$ -- где $F_i$ значет рационалне функције својих аргументов:
\begin{equation}
\sigma_1(x)=F_1(\varphi_1(x)\cdots\varphi_n(x));\quad \cdots;\quad
\sigma_n(x)=F_n(\varphi_1(x)\cdots\varphi_n(x));
\tag{5}
\end{equation}
доколе обратно $\varphi_1(x)\cdots\varphi_n(x)$, и следователно такоже $\varphi(x)=\mu_1\varphi_1(x)+\cdots+\mu_n\varphi_n(x)$, алгебраично зависет од $\sigma_1(x)\cdots\sigma_n(x)$ посредством уравненьа, неразложимого в односу к $\Omega(\sigma_1(x)\cdots\sigma_n(x))$:
\begin{equation}
G(\varphi(x);\sigma_1(x)\cdots\sigma_n(x))=0,
\tag{6}
\end{equation}
кторо знову јест идентичност в $x$.

Посредством (5) та идентичност (6) преходи в
\begin{equation}
G(\varphi(x);F_1(\varphi_1\cdots\varphi_n);\cdots;F_n(\varphi_1\cdots\varphi_n))
=H(\varphi_1(x)\cdots\varphi_n(x))=0;
\tag{7}
\end{equation}
и заради алгебраичној независности $\varphi_1(x)\cdots\varphi_n(x)$ (7) јест исполньено не само идентично в $x$, але такоже в $\varphi_i$; т. ј. идентично в независных параметрах $\lambda$ важи
\begin{equation}
H(\lambda_1\cdots\lambda_n)=G(\mu_1\lambda_1+\cdots+\mu_n\lambda_n;F_1(\lambda_1\cdots\lambda_n);\cdots;F_n(\lambda_1\cdots\lambda_n))=0.
\tag{8}
\end{equation}
Из идентичности (8) следује, же уравненье
\begin{equation}
f(x)=x^n-F_1(\lambda_1\cdots\lambda_n)x^{n-1}+F_2(\lambda_1\cdots\lambda_n)x^{n-2}\cdots\pm F_n(\lambda_1\cdots\lambda_n)=0
\tag{9}
\end{equation}
{\itshape в односу к} $\Omega(\lambda_1\cdots\lambda_n)$ {\itshape имаје групу} $\Gamma$. Бо та идентичност показује, же рационално знано $\lambda=\mu_1\lambda_1+\cdots+\mu_n\lambda_n$ јест функција коренов принадлежеча $\Gamma$; але далье не може быти рационално знана никака функција принадлежеча подгрупе $\Gamma$, како показује субституција $\lambda_i=\varphi_i(x)$; а при тој самој субституцији $\lambda$ такоже јест различна од всих својих конјугатов. Ако далье $\Omega^*$ јест либовольно числово полье садржече $\Omega$, тогда $f(x)$ {\itshape имаје групу} $\Gamma$ {\itshape в односу к} $\Omega^*(\lambda_1\cdots\lambda_n)$; бо по ч. 2 група не јест редуцирана при субституцији $\lambda_i=\varphi_i(x)$ ни после адјункције $\Omega^*$.

Оставаје доказати, же в смыслу модификованом в ч. 1, $f(x)$ стварно пробегаје целост всих $g_\Gamma$. За то заметимо, же по (5) функције $F_i(\lambda)$ сут алгебраично независне функције својих аргументов $\lambda$. Систем уравнень
\begin{equation}
F_1(\lambda_1\cdots\lambda_n)=-a_1;\quad F_2(\lambda_1\cdots\lambda_n)=a_2;\quad \cdots;\quad F_n(\lambda_1\cdots\lambda_n)=\pm a_n
\tag{10}
\end{equation}
има зато решеньа за всакы либовольны систем вредностиј $a_1\cdots a_n$ -- с изкльученьем јешче означимых сингуларных. За всакы систем вредностиј $a_1\cdots a_n$, кторы одповедаје $g_\Gamma$ в односу к $\Omega^*$, одповедајуче $\lambda_i$ како природне ирационалности принадлежече групе\footnote{Срав. такоже ниже!} ставајут величинами из $\Omega^*$. Тогда, ако $\lambda_1\cdots\lambda_n$ пробегајут все системы вредностиј, $g(x)$ пробегаје целост всих уравнень в модификованом смыслу; и та целост садржаје целост $g_\Gamma$, кторым одповедајут вредности из $\Omega^*$. Понеже обратно, по Хилбертовом теорему о неразложимости, вредностјам $\lambda$ из $\Omega^*$ «в обче» такоже одповедајут уравненьа $g_\Gamma$, доставаје се теды в смыслу модификованом в ч. 1 чрез параметрно представјенье (9) стварно целост всих $g_\Gamma$ в односу к $\Omega^*$.

Оставаје јешче определити сингуларне системы вредностиј, за кторе параметрно представјенье не успеваје. Нехај функције минималној базы, раздельене на числитель и именователь, будут дане чрез
\[
\varphi_i(x)=\frac{\psi_i(x)}{\chi_i(x)}.
\]
Ако то поставимо в идентичности (5), нехај ты без скраченьа преходет в
\begin{equation}
\sigma_k(x)=\frac{G_k(x)}{H_k(x)}\quad\text{или}\quad H_k(x)\cdot\sigma_k(x)=G_k(x),
\tag{11}
\end{equation}
где продукт всих $H_k(x)$ јест делимы продуктом всих именовательев $\chi_i(x)$. За все системы коренов $\alpha_1\cdots\alpha_n$, за кторе продукт
\[
H_1(\alpha)\cdot H_2(\alpha)\cdots H_n(\alpha)\neq0
\]
и следователно никакы именователь $\chi_i(\alpha)$ не изчезаје, можно в (11) делити чрез $H_k(\alpha)$ и достати
\begin{equation}
\sigma_k(\alpha)=\frac{G_k(\alpha)}{H_k(\alpha)}=F_k(\varphi_1(\alpha);\cdots;\varphi_n(\alpha)),
\tag{12}
\end{equation}
при чем $\varphi_i(\alpha)$ заради неизчезаньа именовательев пријмајут конечне вредности -- и при уравненьах $g_\Gamma$ ставајут величинами из $\Omega^*$. (12) представльа систем решень (10), скоро $\alpha_1\cdots\alpha_n$ сут определьене тако, же $a_k=(-1)^k\sigma_k(\alpha)$.\footnote{Решенье уравненьа тым јест сведено на решенье система независных функциј:
\[
\varphi_1(\alpha_1\cdots\alpha_n)=\lambda_1;\quad \cdots;\quad \varphi_n(\alpha_1\cdots\alpha_n)=\lambda_n.
\]}

{\itshape Сингуларне} могут теды быти само таке системы коефициентов $a_k=(-1)^k\sigma_k(\alpha)$, за кторе
\[
H_1(\alpha)\cdot H_2(\alpha)\cdots H_n(\alpha)=0;
\]
за кторе теды (11), в место представјеньа, преходи в $0=0$. Тогда јест потребно особно изследованье, чи меджу тако дефиноваными сингуларными $a_1\cdots a_n$ сут такоже таке, кторым одповедајут уравненьа $g_\Gamma$, и за кторе числове польа.\footnote{Напр., како Ф. Сеиделманн показал на наведеном месте, само при алтернативној групе уравнень 3-го и 4-го степена наступајут така уравненьа -- одповедајуча сингуларным системам вредностиј -- и то само за така числове польа $\Omega^*$, кторе садржајут полье третјих коренов јединицы. В всаком случају можно дати просто дополнително представјенье.} Соберајучи, доставаје се

{\itshape Теорем. Ако Лагранжев родовы домен принадлежечи групе $\Gamma$ имаје минималну базу}\footnote{Из обстојаньа једнеј минималној базы одмах следује обстојанье либовольно многих, кторе можно извести из изходној чрез обратиме рационалне трансформације.} {\itshape с коефициентами из $\Omega$, тогда за всако числово полье $\Omega^*$ садржече $\Omega$ целост уравнень с заданој групој $\Gamma$ в односу к $\Omega^*$ можно рационално конструовати чрез параметрно представјенье (2) -- евентуално с изкльученьем тых, кторе сут сингуларне в односу к параметрному представјеньу и кторе можно дати особно. Параметры при том имајут пробегати все вредности из $\Omega^*$, кторе не вызывајут редукцију групы. Конструкција јест можна за всако либовольно числово полье, ако коефициентна област $\Omega$ минималној базы јест равна польу $\mathfrak R$ всих рационалных чисел.}

Понеже по Хилбертовом теорему о неразложимости параметрна многовидност не понижаје се выкльученьем тых вредностиј из $\Omega^*$, кторе вызывајут редукцију групы, важи јешче: {\itshape из обстојаньа минималној базы следује, же многовидност уравнень с групој $\Gamma$ в односу к $\Omega^*$ јест равна многовидности уравнень без афекта; многовидност не понижаје се афектом.}

\begin{center}
§ 2.\\[0.75\baselineskip]
\textbf{Примененье к специалным уравненьам.}
\end{center}

Тепер показујемо обстојанье минималној базы за все групы, кторе наступајут при уравненьах 3-го и 4-го степена; и за то најпрво вполно обче показујемо, же за уравненьа $n$-го степена проблем можно редуцирати на таковы од $n-2$ неизвестных. Прва редукција одповедаје линеарној трансформацији Чирнхауса за одстраньенье другого највышего члена. Беремо како прву базову функцију $\varphi_1(x)=\sigma_1(x)$; и далье замечајемо, же $\mathfrak L_\Gamma$ можно рационално извести из всих {\itshape хомогенных} форм од $x_1\cdots x_n$, кторе допушчају $\Gamma$. Нехај $p(x_1\cdots x_n)$ буде така форма; тогда такоже
\[
q(x)=p\left(x_1-\frac{\sigma_1(x)}{n};\cdots;x_n-\frac{\sigma_1(x)}{n}\right)=p\left(x_i-\frac{\sigma_1(x)}{n}\right)=p(y_i)
\]
принадлежи к $\mathfrak L_\Gamma$, понеже она допушча $\Gamma$, и надто имајемо
\[
p(x)\equiv p\left(x_i-\frac{\sigma_1(x)}{n}\right)\operatorname{mod}\sigma_1(x)
\]
или
\[
p(x)=q(x)+\sigma_1(x)\cdot p_1(x)
\]
с $p_1(x)$ принадлежечим к $\mathfrak L_\Gamma$ и меншего степена неж $p(x)$. Продолженье поступка показује, же все хомогене формы из $\mathfrak L_\Gamma$ -- и следователно в обче все рационалне функције из $\mathfrak L_\Gamma$ -- можно рационално изразити чрез $\varphi_1(x)=\sigma_1(x)$ и чрез хомогене формы
\[
q(x)=p\left(x_i-\frac{\sigma_1(x)}{n}\right)=p(y_i).
\]
Меджу теми величинами $y_i$ обстојит вшак линеарна релација $\sum y_i=0$; $q(x)$ {\itshape зависет теды само од $(n-1)$ аргументов.}

Друга редукција опираје се на том, же $q(x)$ сут хомогене формы својих аргументов. Беремо како другу базову функцију
\[
\begin{aligned}
\varphi_2(x)&=\frac{\tau_3(x)}{\tau_2(x)},\\
\tau_3(x)&=\sigma_3\left(x_i-\frac{\sigma_1(x)}{n}\right)=\sigma_3(y_i),\\
\tau_2(x)&=\sigma_2\left(x_i-\frac{\sigma_1(x)}{n}\right)=\sigma_2(y_i).
\end{aligned}
\]
$\varphi_2(x)$ јест теды хомогенно-дробна функција првого степена од $(n-1)$ аргументов, именне од $y_i$ с $\sum y_i=0$. Посредством $\varphi_2(x)$ можно все $q(x)$ рационално изразити чрез функције, такоже принадлежече к $\mathfrak L_\Gamma$,
\[
r(x)=\frac{q(x)}{\varphi_2(x)^\lambda}=\frac{q(x)\cdot\tau_2(x)^\lambda}{\tau_3(x)^\lambda}=\frac{p(y_i)\cdot\sigma_2(y_i)^\lambda}{\sigma_3(y_i)^\lambda},
\]
где $\lambda$ означаје степен $q(x)$. $r(x)$ ставаје теды хомогенна нулового степена и зависи само од $(n-2)$ аргументов, од односов $y_1:y_2:\cdots:y_n$ с $\sum y_i=0$. $\mathfrak L_\Gamma$ можно теды рационално изразити чрез
\[
\varphi_1(x)=\sigma_1(x);\quad \varphi_2(x)=\frac{\tau_3(x)}{\tau_2(x)}
\]
и чрез {\itshape полье всих функциј $r(x)$ из $\mathfrak L_\Gamma$, зависечих од тых $(n-2)$ аргументов},\footnote{За уравненье $n$-го степена без другого члена:
\[
f(x)=x^n+a_2x^{n-2}+a_3x^{n-3}+\cdots+a_n=0
\]
тој другој редукцији одповедаје субституција, всегда можна при $a_2\neq0$, $a_3\neq0$,
\[
y=\frac{x\cdot a_2}{a_3},
\]
чрез коју $f(x)$, до фактора, преходи в
\[
g(y)=y^n+\frac{a_2^3}{a_3^2}y^{n-2}+\frac{a_2^3}{a_3^2}y^{n-3}+\frac{a_4\cdot a_2^4}{a_3^4}y^{n-4}+\cdots+a_n\frac{a_2^n}{a_3^n}=0;
\]
т. ј. в уравненье садржече само $n-2$ параметры:
\[
g(y)=y^n+c\cdot y^{n-2}+c\cdot y^{n-3}+c_4y^{n-4}+\cdots+c_n=0.
\]} чрез што желана редукција јест выконана.

За уравненьа 3-го и 4-го степена одтуд особливо доставаје се редукција на функционалне польа од једного односно двух аргументов. За ты вшак всегда обстојит минимална база по теоремах Лурота и Кастелнуова.\footnote{Ј. Лурот: \emph{Beweis eines Satzes über rationale Kurven}, \emph{Math. Ann. 9 (1875)}; Г. Кастелнуово: \emph{Sulla razionalità delle involuzioni piane}, \emph{Math. Ann. 44 (1893)}. Же за польа од болье неж двух неизвестных минимална база не мора обстојати, показал Ф. Енрикес противным примером, \emph{Rend. Acc. Linc., Vol. 21, 21 Jan. 1912}.} В случају једној неизвестној минималну базу можно извести рационалными операцијами, зато она особливо за $\mathfrak L_\Gamma$ садржаје само рационално-числове коефициенты. При двух неизвестных по Кастелнуову могла бы јешче быти база с коефициентами из алгебраичного числового польа $\Omega$; стварно изследованье (срав. Ф. Сеиделманн, наведено место) показује, же и туто за всако $\mathfrak L_\Gamma$ доставаје се база с рационално-числовыми коефициентами. То можно видети скоро без рачуна уже из тога, же за четверну групу можно избрати таку рационално-числову базу, именне
\[
\varphi_1(x)=\sigma_1(x);\quad \varphi_2(x)=\frac{v w}{u};\quad \varphi_3(x)=\frac{u\cdot w}{v};\quad \varphi_4(x)=\frac{u\cdot v}{w},
\]
где
\[
u=x_1+x_2-x_3-x_4;\quad v=x_1-x_2+x_3-x_4;\quad w=x_1-x_2-x_3+x_4,
\]
тако же алтернативна група преходи в цикличну групу трех базовых функциј $\varphi_2,\varphi_3,\varphi_4$; всака група осмого реда преходи в симетричну групу двух из тых функциј, чрез што наставаје редукција на групы уравнень 3-го и 2-го степена. За цикличну групу рационално-числова база уже јест познана од Абела, иако не формулована како така.\footnote{Срав. \emph{Weber Algebra} (мала издаја), § 93, формула (12).} Тако јест доказано: {\itshape целост уравнень 3-го и 4-го степена с заданој групој -- евентуално с изкльученьем тых, кторе сут сингуларне в односу к параметрному представјеньу -- можно рационално конструовати чрез параметрно представјенье за всаку либовольну област рационалности; ньихова многовидност јест равна многовидности уравнень без афекта.}

Стварно изследованье показује (срав. Ф. Сеиделманн, наведено место), же дополнително представјенье потребно јест само за алтернативну групу в односу к областјам рационалности, кторе садржајут полье третјих коренов јединицы, доколе в всих другых случајах параметрно представјенье даје целост без ограниченьа.

Треба јешче заметити, же минимална база такоже јест познана за все {\itshape Абелове групы}.\footnote{Е. Фишер: \emph{Die Isomorphie der Invariantenkörper der endlichen Abelschen Gruppen linearer Transformationen}. \emph{Gött. Nachr. 1915}, и \emph{Zur Theorie der endlichen Abelschen Gruppen}. \emph{Math. Ann. 77 (1915)}.} Туто вшак иде о базу с коефициентами из циклотомичного польа изведеного из груповых характеров. За конструкцију уравнень с заданој групој зато в том случају не можно добити нове резултаты.

\mbox{Göttingen}, јулиј 1916.

\clearpage

\end{document}

